\documentclass[11pt]{article}

\usepackage[utf8]{inputenc}
\usepackage[T1]{fontenc}
\usepackage{amsmath}
\usepackage{amsfonts}
\usepackage{amssymb}
\usepackage{bm}
\usepackage{graphicx}
\usepackage{natbib}
\usepackage{url}
\usepackage{hyperref}
\usepackage[margin=1in]{geometry}
\usepackage{booktabs}
\usepackage{array}
\usepackage{float}
\usepackage{subcaption}
\usepackage{xr}
\externaldocument{supp}

\setlength{\parskip}{6pt}

\begin{document}

\title{Are seasonal epidemics synchronized or only entrained? A US influenza case study}

\author{Yair Daon\thanks{Azrieli Faculty of Medicine, Bar-Ilan
    University, Safed, Israel. Correspondence: yair.daon@gmail.com}}

\date{\today}

\maketitle

\begin{abstract}
Distinguishing entrainment by a shared driver from synchronization is
a recurring problem in the nonlinear dynamics of living systems, from
coupled genetic oscillators to the spread of infectious diseases. We
argue that for seasonal epidemics the two mechanisms are typically
indistinguishable from population-level observations alone. We apply
the Schur complement to the Fisher information matrix of a coupled SIR
model and show that under entrainment, the trajectory sensitivity to
inter-region coupling is absorbed by fitting initial conditions and a
shared national driver. We then analyze 12 seasons of pneumonia and
influenza mortality records from 47 US states. Surprisingly, a
model-free SARIMA shows improved prediction skill for 65\% of
state pairs, while a mechanistic coupled SIR fit identifies coupling
for 80\% of state pairs. However, once a national driver covariate is
included in the analysis, these figures drop to 50\% and 20\%,
respectively. We conclude that inter-region coupling plays no
detectable role for seasonally driven influenza and make a practical
recommendation: pairwise coupling should be dropped from operational
influenza forecast models in accordance with Occam's razor.  More
broadly, we recommend the same Fisher information diagnostic as an
identifiability check for any entrained system where coupling is
suspected.
\end{abstract}

\noindent\textbf{Keywords:} entrainment $|$ synchronization $|$
Fisher information $|$ coupled oscillators $|$ identifiability $|$
seasonal influenza

\section{Introduction}
The spatial coherence of seasonal infectious-disease epidemics across
geographically separated regions has long attracted attention as a
window into the mechanisms that drive their dynamics. In
pre-vaccination England and Wales, hierarchical traveling waves of
measles incidence moved from large to small population centers along
commuting flows~\citep{grenfell2001}; in Niger,
nighttime-lights-derived mobility helps explain seasonal fluctuations
of measles incidence~\citep{bharti2018}. For dengue across the
Americas, the regional coherence of seasonal cycles ``likely reflects
both regional climate variability and human
movement''~\citep{quandelacy2025}.
%% The mobility-network framework for the spatial spread of pathogens
%% has also been studied as a problem in nonlinear dynamics in its own
%% right~\citep{belik2011}.
The unifying interpretation is that observed spatial coherence ---
together with a plausible coupling mechanism --- implies that coupling
plays a major role in the dynamics.

For seasonal influenza in the United States, two mechanisms have been
recognized as driving coherence. Environmental forcing is often
invoked as a driver: absolute humidity is thought to modulate
per-capita transmission rates, entraining state-level
epidemics~\citep{shaman2009, shaman2010}, though the underlying
mechanism remains contested~\citep{marr2019}. Concurrently, inter-state
coupling through human mobility is widely believed to contribute
additionally to the observed coherence of state-level epidemics
\citep{viboud2006, charu2017, dejong2024, zhang2026}. \citet{yang2014}
established the coupled population approach with US commuting flows,
and more recent influenza forecast systems have built on similar
coupling~\citep{pei2018, osthus2021}.

Independently, model-free time-series methods suggest that influenza-like-illness (ILI)
forecasts in the US benefit from including other states as
covariates~\citep{thivierge2025}, with a modest but consistent
improvement in skill. We reproduce this finding for pneumonia and
influenza mortality in Fig.~\ref{fig:sarima_heatmap}: adding one
state's lag-one mortality as a covariate reduces test-period RMSE for
65\% of target-covariate pairs. The consensus picture, then, is that
US seasonal influenza is both entrained by shared environmental
forcing \emph{and} synchronized to varying degrees by inter-state
coupling.

The distinction between ``entrainment and synchronization'' and ``only
entrainment'' is practical: it dictates what a forecast model should
contain. Entrainment by environmental forcing is uncontested; the
question is whether inter-state coupling contributes identifiable
structure to the trajectories. If it does, then coupling belongs in
the model; otherwise, seeking coupling parameters is redundant, and
operational models can drop it. We therefore ask the question of the
title: \textbf{is seasonal influenza synchronized, or is it only
  entrained?}


We approach this question from two angles. First, we formulate a
mechanistic SIR model and apply the Cram\'er-Rao lower bound to
simulations of this model, building on recent identifiability and
uncertainty analyses of compartmental epidemic
models~\citep{edeling2021, kharazmi2021} and extending
information-theoretic identifiability analyses of single-region
epidemic parameters~\citep{parag2022} to inter-regional coupling. We show that
coupling between regions is fundamentally unidentifiable when
epidemics are entrained by common seasonal forcing. Further, we use the structure of the
resulting Fisher information matrix to argue that either (i)
coupling does not affect the disease trajectory, or (ii) its
effect is absorbed by fitting initial conditions. Since mechanistic
forecast routines fit initial conditions from past observations, our
analysis suggests that coupling may, in fact, be rendered
redundant. We proceed to test
this hypothesis on US pneumonia and influenza mortality. We find
maximum-likelihood estimates of coupling from real US data and
show that, overall, surveillance records contain no statistical
evidence of inter-state coupling.

Second, in an independent, model-free line of analysis, we return to
the forecasting framework of Fig.~\ref{fig:sarima_heatmap} and show that
the apparent state-covariate signal carries no information beyond a
national driver covariate. A conditional Granger causality test
\citep{granger1969, geweke1984} then implies forecast models should not
include coupling, only a national driver.


Both lines of analysis support the same reading: the spatial
coherence of US seasonal influenza implies entrainment by shared
environmental forcing rather than synchronization through
inter-state coupling, and the coupling parameter can be dropped from
operational forecast models. The same \textbf{design-stage} check applies well beyond
influenza. Whenever coupled oscillators are driven by common forcing
\citep{pikovsky2001}, e.g.~in neural circuits
\citep{reid2019},
ecological metapopulations
\citep{bjornstad1999, grenfell1998},
and climate networks
\citep{tsonis2004architecture},
our approach can flag parameters that the planned data cannot
support. A redundant parameter does not improve forecasts and should
be removed from forecast models in accordance with Occam's razor.


\begin{figure}[H]
\centering
\includegraphics[width=\linewidth]{../pix/sarima_heatmap.pdf}
\caption{Test-period RMSE change when another state's lagged mortality
  is added as a covariate to a forecast of weekly per-capita
  pneumonia and influenza mortality. Each cell corresponds to a single
  target-covariate pair. Blue indicates the covariate deteriorates the
  forecast, and red indicates the covariate improves the forecast.
  Including another state's mortality as a covariate reduces RMSE for
  65\% of target-covariate pairs. Values clipped to~$\pm
  10\%$. States are sorted by the median effect. Train on 2010--2019,
  test on 2023--2025.}
\label{fig:sarima_heatmap}
\end{figure}




















% ---------------------------------------
% -------------- Results ----------------
% ---------------------------------------
\section{Results}
\subsection{Analysis via Schur complement}
The overall behavior of any entrained coupled system can be understood
via the Schur complement of the Fisher information
matrix\footnote{Readers less familiar with the Fisher information are
advised to review the Methods (specifically Section \ref{subsec:crlb})
for details.}. Partition the $6\times 6$ Fisher information matrix of a
seasonally driven coupled SIR model into a $5\times 5$
nuisance-parameter block and a $1\times 1$ coupling block,
\begin{equation*}
  \mathbf{J} =
  %
  %
  \begin{bmatrix}
    J_{\eta\eta}             & \mathbf{J}_{\eta\theta} \\
    \mathbf{J}_{\theta\eta}  & \mathbf{J}_{\theta\theta}
  \end{bmatrix},
\end{equation*}
where $\eta = (S_1(0), I_1(0), S_2(0), I_2(0), \gamma)$ collects the
four per-region initial conditions and the national driver mixing
weight, and $\theta$ is the coupling parameter. By the Schur
complement, the Fisher information for $\theta$ is 
\begin{equation}\label{eq:schur}
  I(\theta) = J_{\theta\theta} - \mathbf{J}_{\theta\eta}\mathbf{J}_{\eta\eta}^{-1}\mathbf{J}_{\eta\theta}.
\end{equation}


A small Fisher information in Eq.~\eqref{eq:schur} arises for one of
two reasons. Either (i) $J_{\theta\theta}$ is small, so $\theta$ has
no effect on trajectories, or (ii) the $\theta$-sensitivity
$\partial\mu/\partial\theta$ is close to the span of the
nuisance-parameter sensitivities $\{\partial\mu/\partial S_i(0),
\partial\mu/\partial I_i(0), \partial\mu/\partial\gamma\}$ (see SI
Section~\ref{sec:si_schur} for details). We expect the
$\gamma$-sensitivity to matter less here, so under (ii), any effect of
$\theta$ on the disease trajectory is absorbed by fitting initial
conditions. For US influenza, states share strong seasonal forcing
with similar phases (see SI~Fig.~\ref{fig:si_phase_theta}), so we
predict that $\theta$ is either redundant~(i) or unidentifiable when
fitting ICs~(ii); either way, the surveillance record should carry
little information about pairwise coupling.


\subsection{Entrainment alone can make coupling unidentifiable}\label{subsec:crlb_inflated}
Simulations also hint that entrainment is key to understanding the
behavior of entrained coupled epidemic systems. We ask whether the
coupling parameter $\theta$ is even in principle recoverable from
observed incidence generated via a two-region SIR model. In
Fig.~\ref{fig:power} we compare statistical power for identifying
$\theta > 0$ across a grid of $\theta$ and seasonal forcing amplitudes
$\varepsilon$, comparing in-phase forcing ($\varphi_1 =
\varphi_2$) with reverse-phase forcing ($\varphi_1 - \varphi_2 =
\pi$). Under reverse-phase forcing, power reaches $0.95$ over most of
the grid. Thus, a two-region SIR system with distinct phases carries
enough information about $\theta$ to make the coupling parameter
identifiable, at least in principle. Under in-phase forcing, however,
power does not exceed $0.55$ anywhere on the grid. At the red star,
marking a strongly connected flu epidemic ($\theta = 0.05$,
$\varepsilon = 0.517$; see Table~\ref{tab:params}), in-phase power is
$0.39$ against reverse-phase $0.89$. Entrainment severely degrades the
detectability of pairwise coupling, in accordance with the previous
section.

\begin{figure}[H]
\centering
\includegraphics[width=\linewidth]{../pix/power.pdf}
\caption{Statistical power for $H_1:\theta > 0$ at $\alpha = 0.05$,
  with 12 seasons of simulated data. \textbf{Left:} in-phase
  forcing ($\varphi_1 = \varphi_2$). \textbf{Right:} reverse-phase
  forcing ($\varphi_1 - \varphi_2 = \pi$). Initial conditions are
  drawn from the empirical distribution of per-season, per-state
  values (SI~Fig.~\ref{fig:si_ics}). Red star: a strongly driven flu
  epidemic ($\theta = 0.05$, $\varepsilon = 0.517$).}
\label{fig:power}
\end{figure}

Results of the simulations suggest there should be very little benefit
in including another state's lag-one mortality in a mechanistic SIR
fit. But if directional coupling is not the underlying cause for the
reduction in in-sample SIR RMSE, then what is it? Two mechanisms can
produce the apparent advantage of \textbf{coupled} over
\textbf{univariate} we observed in Fig.~\ref{fig:sarima_heatmap} and
Fig.~\ref{fig:theta vs crlb coupled}: (a) detectable coupling, in
which one state's past actually affects the other state's future
beyond its past and shared seasonality; or (b) a shared national
driver, in which one state's lagged value is a noisy proxy for the
current week's flu-season severity. The two mechanisms
make opposite predictions: under (a) we expect a reduction in
in-sample RMSE going from \textbf{national driver} to \textbf{both},
while under (b) we expect no such reduction since the national driver
carries all the information pertaining to the severity. Hence, a
comparison of RMSEs between \textbf{national driver} and
\textbf{both} will allow us to discern (a) from (b).


\subsection{Design and Data}\label{subsec:nested}
We analyze excess pneumonia and influenza mortality from 47 US states
over 12 flu seasons (Fig.~\ref{fig:seasons}).
\begin{figure}[H]
\centering
\includegraphics[width=\linewidth]{../pix/seasons_highlighted.pdf}
\caption{Weekly pneumonia and influenza deaths for four US
  states. Light traces show the full time series; darker segments
  indicate the 20-week flu seasons (starting each November) included
  in the analysis. COVID-19 pandemic years (2020--2022) are
  excluded.}
\label{fig:seasons}
\end{figure}
For each pair of states we perform two parallel analyses --- one
model-free (SARIMA, Section~\ref{subsec:sarimabench}) and one
mechanistic (SIR, Section~\ref{subsec:model}). In each framework we
fit four nested models that differ in whether a coupling term and a
national driver term are included; see Table~\ref{tab:cases}.

\begin{table}[h]
\centering
\renewcommand{\arraystretch}{1.15}
\begin{tabular}{@{}l c c@{}}
\toprule
Model & State coupling & National driver \\
\midrule
\textbf{univariate}                   & $\times$   & $\times$   \\
\textbf{coupled}                        & \checkmark & $\times$   \\
\textbf{national driver}             & $\times$   & \checkmark \\
\textbf{both}   & \checkmark & \checkmark \\
\bottomrule
\end{tabular}
\caption{Four nested  models.}
\label{tab:cases}
\end{table}

We focus on two comparisons:
\begin{enumerate}
\item \textbf{univariate} vs.\ \textbf{coupled}, and
\item \textbf{national driver} vs.\ \textbf{both}.
\end{enumerate}
For example, results for the comparison \textbf{univariate}
vs.~\textbf{coupled} were already presented in
Fig.~\ref{fig:sarima_heatmap}.


\subsection{SIR model: \textbf{univariate} vs.~\textbf{coupled}}\label{subsec:us_real_data}
We consider a coupled SIR model (see Section \ref{subsec:model} and
specifically Eq.~\eqref{eq:foi} for details), restricted to pairwise
coupling only. For each state pair we conduct a comparison of two
competing models (see Section~\ref{subsec:nested}): the
\textbf{coupled} model (with coupling parameter $H_1:\theta>0$)
against a \textbf{univariate} null, where no coupling is assumed
($H_0:\theta=0$). A hypothesis test (see Section \ref{subsec:hyp} for
details) rejects the null $H_0: \theta = 0$ for $80.7\%$ of state
pairs at $\alpha = 0.05$ (see Fig.~\ref{fig:theta vs crlb coupled}),
seemingly affirming the ILI-forecast signal of~\citet{thivierge2025}
on the P\&I mortality endpoint: pairwise coupling looks statistically
significant.

\begin{figure}[H]
\centering
\includegraphics[width=0.75\linewidth]{../pix/theta_vs_crlb_coupled.pdf}
\caption{Scatter plot of $\hat{\theta}$ (x-axis) and standard deviation
  (y-axis) across the 1081 US state pairs. The rejection threshold is
  dotted black. Above the rejection threshold $H_0:\theta=0$ is not
  rejected; below the rejection threshold $H_0:\theta=0$ is rejected
  in favor of the alternative $H_1:\theta > 0$ (see Section
  \ref{subsec:hyp}). The bulk of pairs are below the rejection line
  ($80.7\%$ reject $H_0: \theta = 0$) in accordance with
  \cite{thivierge2025}.}
\label{fig:theta vs crlb coupled}
\end{figure}



\subsection{SIR model: \textbf{national driver} vs.~\textbf{both}}\label{subsec:natavg_results}
Following the reasoning of the previous section we fit each state pair
allowing $\gamma \in [0, 1]$, and conduct the comparison of the null
\textbf{national driver} ($\gamma > 0$, $\theta = 0$) vs. the
alternative \textbf{both} ($\gamma>0$ and $\theta>0$). Our hypothesis
test (see Section \ref{subsec:hyp} for details) rejects the null
\textbf{national driver} for $18.6\%$ of state pairs in favor of
\textbf{both} (compared to $80.7\%$ rejecting
\textbf{univariate} in favor of \textbf{coupled}); see
Fig.~\ref{fig:theta vs crlb both}.
\begin{figure}[H]
\centering
\includegraphics[width=0.75\linewidth]{../pix/theta_vs_crlb_both.pdf}
\caption{Scatter plot of $\hat{\theta}$ (x-axis) and standard deviation
  (y-axis) across the 1081 US state pairs. The rejection threshold is
  dotted black. Above the rejection threshold $H_0:\theta=0$ is not
  rejected; below the rejection threshold $H_0:\theta=0$ is rejected
  in favor of the alternative $H_1:\theta > 0$ (see Section
  \ref{subsec:hyp}). The bulk of pairs sit above the rejection
  threshold ($18.6\%$ reject $H_0: \theta = 0$). Compare to
  $80.7\%$ in Fig.~\ref{fig:theta vs crlb coupled}.}
\label{fig:theta vs crlb both}
\end{figure}


\subsection{SARIMA: \textbf{national driver} vs.~\textbf{both}}\label{subsec:sarima_alpha}
We now revisit the model-free autoregressive framework of
Fig.~\ref{fig:sarima_heatmap}. Again, we compare a null
\textbf{national driver} to the alternative \textbf{both}. We note
that this comparison is in fact a conditional Granger causality
test~\citep{granger1969, geweke1984}. We had already established in
Fig.~\ref{fig:sarima_heatmap} that $y_{B,t-1}$ improves the forecast
of $y_{A,t}$, so by Reichenbach's common cause principle and temporal
ordering either $y_{B,t-1}\to y_{A,t}$, or there is a common cause. If
the forecast skill using $y_{B,t-1}$ and $\bar{y}_{t-1}$ is no better
than that of $\bar{y}_{t-1}$, then $\bar{y}_{t-1}$ is the common cause
and $Y_B \leftarrow \bar Y \to Y_A$~\citep{pearl2009}.  In this case,
the apparent $Y_B \to Y_A$ relationship is confounded by the national
driver severity rather than reflecting directional inter-state
coupling.

\begin{figure}[H]
\centering
\includegraphics[width=\linewidth]{../pix/sarima_density.pdf}
\caption{Distribution of relative test-period RMSE change
  \eqref{eq:sarima_delta} when the covariate state's lag-one mortality
  $y_{B,t-1}$ is added as a second covariate to a SARIMA that already
  includes the lag-one national driver $\bar y_{t-1}$
  across target-covariate pairs. Negative values imply that adding the
  covariate state improves the forecast. The distribution is nearly
  symmetric around zero (median $+0.01\%$, dashed red line), showing
  that the covariate state carries no information beyond the
  national driver.}
\label{fig:sarima_density}
\end{figure}

Fig.~\ref{fig:sarima_density} discerns the two options: the
information in $y_{B,t-1}$ is merely a proxy for the national driver,
which the national driver captures cleanly. Had coupling
been an important factor, $y_{B,t-1}$ would have carried information
not contained in $\bar y_{t-1}$ and the joint model would beat the
national driver-only model. We see no such effect. Conditional Granger
causality then implies the data are consistent with a common-cause
graphical structure $Y_B \leftarrow \bar Y \to Y_A$, in which the
national driver confounds the apparent directed-coupling
structure $Y_B \to Y_A$.


\section{Discussion}
Our complementary analyses --- a mechanistic SIR fit and a model-free
conditional Granger test --- arrive at an identical conclusion:
granted the established role of shared environmental forcing in
driving US seasonal influenza, the additional contribution of
inter-state coupling is practically unidentifiable from observed
disease trajectories once a shared national driver is accounted for,
and \textbf{the data are consistent with entrainment
  alone}. Inter-state coupling can therefore be dropped from
operational forecast models without loss of identifiable predictive
structure.

The naive coupled SIR fit rejects $H_0: \theta = 0$ for a large
fraction of state pairs, in apparent contradiction with our Fisher
information Schur complement prediction that $\theta$ should be
unidentifiable under shared seasonal forcing. The conditional Granger
test on SARIMA locates the source: the state-covariate signal is a
proxy for national driver, not a signature of directional
inter-state coupling. Extending the mechanistic FOI with a national
driver dissolves the apparent coupling signal from within the SIR
framework: once the shared driver is explicit in the model,
$\hat\theta$ falls below the rejection threshold for most state pairs,
recovering the theoretical prediction. In this sense $\theta$ is not
so much identifiable as it is a receptacle for shared driving that has
nowhere else to go.

Recently, \citet{thivierge2025} asked whether ILI forecasting in US
states benefits from including other-state or US-weighted-average ILI
as covariates in various autoregressions. They found consistent
improvements by both RMSE and weighted interval score, and importantly
observed that other-state and US-weighted-average covariates produce
\emph{similar} gains. In our SARIMA framework, the national driver
covariate alone produces a median RMSE reduction of $7.0\%$ and
reduces RMSE in $95\%$ of target-covariate pairs, while individual
states contribute essentially nothing beyond the national driver. Our
findings therefore complement those of \citet{thivierge2025} at the
mortality endpoint: the other-state gains they observed are explained
by a shared national driver rather than inter-state coupling, and our
$\gamma$-extended SIR fit provides the mechanistic reason --- coupling
and a shared driver are exchangeable in fitting the observed
trajectory, and only one of them is required.

Our analysis also carries implications for influenza forecasting in
tropical regions~\citep{tsang2024}. In tropical regions, weaker and
less coherent seasonal forcing may produce reverse-phase
outbreaks~\citep{tamerius2013}, which our analysis shows may give rise
to the phase asymmetry needed for identifiability. Inter-state
coupling may therefore be identifiable in tropical regimes --- and
contribute genuine forecast skill --- where it cannot in the entrained
temperate regime. Our Cram\'er-Rao analysis offers a precise way to
characterize this distinction.


\subsection{Beyond Influenza}
The identifiability-under-entrainment diagnostic we apply here is not
specific to influenza or to inter-regional coupling. Contemporary spatial and structured
epidemic models routinely carry coupling matrices, age-specific
contact rates, and more. Because the Cram\'er-Rao lower bound bounds
the variance of every unbiased estimator and is computable from the
model and a hypothesized parameter value alone, the same design-stage
check can be applied to any such model before any inference pipeline
is chosen: a parameter whose effect on the trajectory can be
reproduced by fitting other parameters does not improve forecasts and
should be removed from the model in accordance with Occam's razor.
Recent work has begun to characterize uncertainty and identifiability
in compartmental epidemic models~\citep{edeling2021, kharazmi2021},
identifiability constraints in network-based epidemic
models~\citep{kiss2023parameter, dilauro2020, kiss2024towards}, and
the informativeness of single-region surveillance data for
reproduction-number inference via Fisher
information~\citep{parag2022}, building on earlier results that
transmission rate and network structure are fundamentally confounded
\citep{welch2011}, together providing a parallel design-stage program
for network parameters. The mathematical structure underlying our
particular finding---oscillators driven by common periodic
forcing---also arises outside epidemiology, in neural circuits
receiving shared input
~\citep{reid2019, pikovsky2001}
and in metapopulations driven by environmental
cycles
~\citep{bjornstad1999};
the same entrainment-driven redundancy is plausible there, and the
same design-stage diagnostic applies. Conversely, data-driven methods
for inferring network dynamics from time series~\citep{gao2022} are
likely to succeed where common periodic forcing is absent or weak,
suggesting a natural complementarity with the design-stage diagnostic
proposed here.

\subsection{Limitations}
Several limitations warrant discussion. First, our model assumes a
single influenza strain with uniform transmission parameters across
all seasons and regions. In reality, $R_0$ varies substantially across
influenza subtypes~\citep{biggerstaff2014}. Just like for
coupling, this variability could be absorbed by the initial
conditions.

Second, our observation model is a Gaussian pseudo-likelihood matched
to a negative-binomial mean--variance relationship with dispersion
parameter $k = 10$. This accommodates the overdispersion typical of
surveillance counts and remains valid for the continuous,
sometimes-negative excess-mortality observations produced by baseline
subtraction.

Third, we focused on a two-region model with symmetric coupling.
Extensions to asymmetric coupling or larger networks will likely
reveal additional challenges.

\section{Materials and Methods}

\subsection{Seasonal ARIMA Forecast}\label{subsec:sarimabench}
We fit a seasonal ARIMA model to each state's weekly per-capita
pneumonia and influenza mortality, with and without covariates drawn
from other states. For a target state $A$, the model of a $d$-times
differenced time series is
\begin{equation}\label{eq:sarima_model}
  y_{A,t}  = a_0
    + \sum_{k=1}^{p} \phi_k\, y_{A, t-k}
    + \sum_{j=1}^{q} \psi_j\, \epsilon_{A, t-j}
    + \Phi\, y_{A, t-52}
    + \boldsymbol{\beta}^{\top} \mathbf{x}_{t-1}
    + \epsilon_{A,t}%,
\end{equation}
which is a SARIMA$(p,d,q)(1,0,0)_{52}$ model with a one-week-lagged
exogenous regressor vector $\mathbf{x}_{t-1}$.
The non-seasonal orders $(p, d, q)$ and the coefficients are chosen on
the training data by minimizing AIC~\citep{hyndman2008}.

For each ordered pair of states $(A, B)$ with $B \neq A$, we
fit~\eqref{eq:sarima_model} under four configurations of the exogenous
regressor $\mathbf{x}_{t-1}$; see Table~\ref{tab:SARIMAs}. This
decomposition is the SARIMA instantiation of the four cases in
Table~\ref{tab:cases}.

\begin{table}[h]
\centering
\renewcommand{\arraystretch}{1.15}
\begin{tabular}{@{}l c  c@{}}
\toprule
Model & $\mathbf{x}_{t-1}=$ \\
\midrule
\textbf{univariate}                   & $\emptyset$     \\
\textbf{coupled}                        & $(y_{B, t-1})$   \\
\textbf{national driver}             & $(\bar y_{t-1})$ \\
\textbf{both}   & $(y_{B, t-1},\bar y_{t-1})$ \\
\bottomrule
\end{tabular}
\caption{The four SARIMA instances of the cases in
  Table~\ref{tab:cases}.}
\label{tab:SARIMAs}
\end{table}

Here $\bar y_{t}$ is the population-weighted per-capita mortality
across all states other than the pair $(A, B)$ under test,
\begin{equation*}
\bar y_{t} \;=\; \frac{\sum_{S \notin \{A, B\}} y_{S, t}\, N_S}{\sum_{S \notin \{A, B\}} N_S},
\end{equation*}
computed identically to the national driver of the mechanistic model
(Section~\ref{subsec:model}). In $\mathbf{x}_{t-1}$ we use it at lag
one. We train on 2010--2019 and test on 2023--2025, and report the
relative change in test-period RMSE,
\begin{equation}\label{eq:sarima_delta}
  \Delta \;=\; \frac{\text{RMSE}_{\mathrm{alt}} - \text{RMSE}_{\mathrm{ref}}}{\text{RMSE}_{\mathrm{ref}}},
\end{equation}
so $\Delta < 0$ means the alternative model improves the forecast over
the reference.

\subsection{Mechanistic Model}\label{subsec:model}
Consider two interconnected regions with populations $N_1$ and $N_2$.
Let $S_i(t)$, $I_i(t)$, $R_i(t)$ denote the number of susceptible,
infected, and recovered individuals in region $i$ at time $t$, with
$S_i + I_i + R_i = N_i$. Regional coupling is governed by the contact
matrix
%
%
$$
\mathbf{C} =
\begin{bmatrix}
  1 - \theta  &     \theta \\
      \theta  & 1 - \theta
\end{bmatrix}
$$
%
%
where $\theta \in [0, 1]$ is the coupling parameter. The
transmission rate is a standard seasonal forcing ($\varepsilon$ is forcing
amplitude, $\varphi_i$ is a phase offset):
%
%
$$
\beta_i(t) = \beta_0(1 + \varepsilon\sin(2\pi t + \varphi_i)),
$$
%
%
where $t$ is measured in years. We let $I_{\text{nat}}$ and
$N_{\text{nat}}$ denote the total number of nationally infected
individuals and the national population, respectively. The
force of infection in region $i$ is~\citep{edlund2011}
\begin{equation}\label{eq:foi}
\lambda_i(t) = \beta_i(t)\left[(1-\gamma)\frac{\sum_j C_{ij} I_j(t)}{\sum_j C_{ij} N_j}
+ \gamma \frac{I_{\text{nat}}(t)}{N_{\text{nat}}}\right],
\end{equation}
where $\gamma \in [0, 1]$. We estimate the driver from aggregate P\&I
excess mortality across all states \emph{other than the simulated
pair} as $I_{\text{nat}}(t) = Y_{\text{nat}}(t)/\rho$, where
$Y_{\text{nat}}(t) = \sum_{S \notin \{A, B\}} y_{S,t}$ is total weekly
excess mortality summed across all included states except the pair,
and $N_{\text{nat}} = \sum_{S \notin \{A, B\}} N_S$ is the
corresponding total population.

The expected incidence is $\mu_i(t) = S_i(t)[1 -
  \exp(-\lambda_i(t))]$, and the recovery probability per time step is
$\exp(-\nu)$. Time steps are denoted $\Delta t$, and the resulting
discrete-time SIR dynamics are:
\begin{align}\begin{split}\label{eq:discrete_sir}
S_i(t+\Delta t) &= S_i(t) \exp(-\lambda_i(t)) \\
I_i(t+ \Delta t) &= I_i(t) \exp(-\nu) + \mu_i(t)
\end{split}\end{align}
We use SIR rather than SIRS dynamics because we conduct inference for
each influenza season separately. Model parameters are calibrated to
seasonal influenza following~\citep{shaman2010, rolfes2018}; see
Table~\ref{tab:params}.

\begin{table}[h]
\centering
\begin{tabular}{llll}
\toprule
Parameter & Value & Description & Source / calc.\\
\midrule
$R^{\text{max}}_0$ & 3.52 & $R_0$ at peak season & \citep{shaman2010} \\
$R^{\text{min}}_0$ & 1.12 & $R_0$ during the summer & \citep{shaman2010} \\
$D$ & 3.24 & Mean infection period (days) & \citep{shaman2010} \\
$\nu$ & $2.16$ & Weekly recovery rate & $7 / D$ \\
$\beta_{\text{max}}$ & $3.11$ & Peak-season transmission & $R_0^{\text{max}}(1 - e^{-\nu})$\\
$\beta_{\text{min}}$ & $0.99$ & Low-season transmission & $R_0^{\text{min}}(1 - e^{-\nu})$ \\
$\beta_0$ & $2.05$ & Mean transmission rate & $(\beta_{\text{min}} + \beta_{\text{max}})/2$ \\
$\varepsilon$ & $0.517$ & Seasonal amplitude & $\beta_{\text{max}} / \beta_0 - 1$\\
$\rho$ & $5.36 \times 10^{-4}$ & Infection fatality rate & \citep{rolfes2018}\\
$\varphi$ & & Phase of seasonal forcing & per state \\
\bottomrule
\end{tabular}
\caption{Model parameters calibrated to seasonal influenza.}
%% with $\beta(t) = \beta_0(1+\varepsilon \sin(2\pi t + \varphi))$.}
\label{tab:params}
\end{table}


\subsubsection{Observations and Likelihood}\label{subsub:noise}
Each infection is reported independently with probability $\rho$;
however, surveillance counts exhibit overdispersion beyond the
binomial baseline due to reporting heterogeneity and other noise
sources not captured by the mean-field SIR dynamics. We model
observations as having conditional mean $\rho\mu_i(t)$ and conditional
variance matching a negative-binomial distribution with dispersion
parameter $k$:
%
%
$$
Y_i(t) = \rho\mu_i(t) + \epsilon_i(t)
$$
%
%
with
%
%
$$
\epsilon_i(t) \sim \mathcal{N}\left(0,\; \rho\mu_i(t) +
\frac{(\rho\mu_i(t))^2}{k}\right).
$$
%
%
We fix $k = 10$ throughout.

The unknown parameters for a single influenza season are
$\boldsymbol{\phi} = (S_1(0), I_1(0), S_2(0), I_2(0), \gamma, \theta)$, giving
the log-likelihood:
\begin{equation}\label{eq:loglik}
\ell(\boldsymbol{\phi}) = -\frac{1}{2}\sum_{t=0}^{T-1}\sum_{i=1}^{2}
\left[\log(2\pi \sigma_i^2(t)) + \frac{(Y_i(t) -
    \rho\mu_i(t))^2}{\sigma_i^2(t)}\right]
\end{equation}
where $\sigma_i^2(t) = \rho\mu_i(t) + (\rho\mu_i(t))^2/k$.

\subsubsection{Derivatives}\label{subsub:sens}
We require partial derivatives of model outputs with respect to
parameters in order to find the Fisher information and for maximizing
the likelihood (details below for both). Let $\phi \in \{S_1(0),
I_1(0), S_2(0), I_2(0), \gamma, \theta\}$ denote a parameter ($\phi =
\phi_k$ for $k = 1, \dots, 6$). It is convenient to write the
FOI~\eqref{eq:foi} as $\lambda_i = \beta_i(t)\bigl[(1-\gamma)\,L_i +
  \gamma\,\tilde{N}_i\bigr]$ where $L_i(t) =
(\mathbf{C}\mathbf{I})_i/(\mathbf{C}\mathbf{N})_i$ is the local,
contact-weighted per-capita infected fraction and $\tilde{N}_i(t) =
I_{\text{nat}}(t)/N_{\text{nat}}$ is the exogenous national driver
(independent of $S_i, I_i$ and of $\gamma, \theta$).  For any $\phi
\in \{S_1(0), I_1(0), S_2(0), I_2(0), \theta\}$,
\begin{equation}\label{eq:sens_foi}
\frac{\partial \lambda_i}{\partial \phi} = \beta_i(t)\,(1-\gamma)\left[
\frac{1}{(\mathbf{C}\mathbf{N})_i}\left(
  \frac{\partial \mathbf{C}}{\partial \phi} \mathbf{I}
  + \mathbf{C} \frac{\partial \mathbf{I}}{\partial \phi}\right)_i
- \frac{L_i}{(\mathbf{C}\mathbf{N})_i}
  \left(\frac{\partial \mathbf{C}}{\partial \phi} \mathbf{N}\right)_i
\right];
\end{equation}
for $\phi = \gamma$, the contact matrix has no $\gamma$-dependence and
$\tilde{N}_i$ contributes an additive term,
\begin{equation}\label{eq:sens_foi_alpha}
\frac{\partial \lambda_i}{\partial \gamma} = \beta_i(t)\left[
\tilde{N}_i - L_i
+ \frac{(1-\gamma)}{(\mathbf{C}\mathbf{N})_i}
\left(\mathbf{C}\,\frac{\partial \mathbf{I}}{\partial \gamma}\right)_i
\right].
\end{equation}
The contact matrix derivative is
\begin{align*}
\partial \mathbf{C}/\partial \theta =
\begin{bmatrix}
  -1 &  1 \\
  \phantom{-}1 & -1
\end{bmatrix},\\
\qquad
\partial \mathbf{C}/\partial \phi = \mathbf{0} \text{ for } \phi \neq \theta.
\end{align*}
Partial derivatives of the expected incidence are found from the
product rule:
\begin{equation}\label{eq:sens_mu}
  \frac{\partial \mu_i}{\partial \phi} =
  %
  %
  \frac{\partial S_i}{\partial \phi} [1 - e^{-\lambda_i}] + S_i
  e^{-\lambda_i} \frac{\partial \lambda_i}{\partial \phi}.
\end{equation}
Partials of the system state propagate via:
\begin{align}\begin{split}\label{eq:sens}
\frac{\partial S_i(t+\Delta t)}{\partial \phi} &= e^{-\lambda_i}
\frac{\partial S_i}{\partial \phi} - S_i e^{-\lambda_i}
\frac{\partial \lambda_i}{\partial \phi} \\
\frac{\partial I_i(t+\Delta t)}{\partial \phi} &= e^{-\nu}
\frac{\partial I_i}{\partial \phi} +
\frac{\partial \mu_i}{\partial \phi}
\end{split}\end{align}
with initial conditions $\partial S_i(0)/\partial S_j(0) =
\delta_{ij}$, $\partial I_i(0)/\partial I_j(0) = \delta_{ij}$, where
$\delta_{ij}$ is the Kronecker delta. We solve for these partial
derivatives jointly with the state equations. These
sensitivity-equation methods have been reviewed
by~\citep{wu2013sensitivity} and recently extended
in~\citep{mester2022}.

\subsection{Fisher Information}\label{subsec:crlb}
%% For a single epidemic season, the Cram\'er-Rao Lower Bound (CRLB)
%% states that $\text{Var}(\hat\phi_k) \geq [\mathbf{J}^{-1}]_{kk}$,
We would like to calculate the Fisher information $I(\theta)$ for two
reasons. First, under standard regularity conditions, the
maximum-likelihood estimator $\hat\theta_{\text{MLE}} \sim
\mathcal{N}(\theta, I(\theta)^{-1})$
(asymptotically)~\citep{wilke2020}. Second, the Fisher information
gives rise to the Cram\'er-Rao lower bound which bounds the variance
of any unbiased estimator from below
\begin{equation}\label{eq:crlb}
  \text{Var}(\hat{\theta}) \geq I(\theta)^{-1}.%[\mathbf{J}^{-1}]_{\theta\theta}.
\end{equation}

We let $\mathbf{J}$ denote the $6 \times 6$ Fisher information matrix with
entries $J_{kl} =
\mathbb{E}[(\partial\ell/\partial\phi_k)(\partial\ell/\partial\phi_l)]$,
indexed by $\boldsymbol{\phi} = (S_1(0), I_1(0), S_2(0), I_2(0),
\gamma, \theta)$.  Let $r_i(t) := Y_i(t) - \rho\mu_i(t)$ denote the
residual and
%
%
\begin{equation*}
a_i(t) := \frac{\partial \sigma_i^2(t)}{\partial \mu_i(t)} = \rho + \frac{2\rho^2 \mu_i(t)}{k}
\end{equation*}
%
%
denote the derivative of the variance function with respect to the
mean. One can verify that if we define
\begin{equation*}
A_i(t) := -\frac{a_i(t)}{2\sigma_i^2(t)} + \frac{\rho\, r_i(t)}{\sigma_i^2(t)}
+ \frac{r_i^2(t)\,a_i(t)}{2\sigma_i^4(t)},
\end{equation*}
then differentiating the log-likelihood~\eqref{eq:loglik} yields
%
%
\begin{equation}\label{eq:grad}
\partial \ell / \partial \phi = \sum_{t,i} A_i(t)\,\partial
\mu_i(t)/\partial \phi.
\end{equation}
%
%
Observations are independent across $(i,t)$, and we can easily verify
$\mathbb{E}[A_i(t)] = 0$. We obtain, using the central moments of the
Gaussian $r_i(t) = Y_i(t) - \rho\mu_i(t)$, namely
%% $\mathbb{E}[r_i(t)] = \mathbb{E}[r_i^3(t)] = 0$,
%% $\mathbb{E}[r_i^2(t)] = \sigma_i^2(t)$, and
$\mathbb{E}[r_i^4(t)] = 3\sigma_i^4(t)$, together with~\eqref{eq:grad}:
\begin{align}\label{eq:fim_scalar}
  \begin{split}
    J_{kl} &= \mathbb{E}\left [\frac{\partial \ell}{\partial \phi_k}\frac{\partial \ell}{\partial \phi_l}\right ] \\
    %
    %
    %
    &=\sum_{t,i} \mathbb{E}[A_i^2(t)] \frac{\partial \mu_i(t)}{\partial \phi_k}  \frac{\partial \mu_i(t)}{\partial \phi_l} \\
    %
    %
    %
    &= \sum_{t=0}^{T-1}\sum_{i=1}^{2}\left( \frac{\rho^2}{\sigma_i^2(t)} +
    \frac{a_i^2(t)}{2\sigma_i^4(t)}\right) \frac{\partial
      \mu_i(t)}{\partial \phi_k} \frac{\partial \mu_i(t)}{\partial
      \phi_l}
  \end{split}
\end{align}
%
%
In matrix form, \eqref{eq:fim_scalar} can be written as 
\begin{align}\label{eq:fim}
  \begin{split}
  \mathbf{J} &= \mathbf{G}^T \mathbf{W} \mathbf{G} \\
  &\text{ where } \\
  %
  %
  \mathbf{G} &= \begin{pmatrix}
    \dfrac{\partial \mu_1(0)}{\partial S_1(0)} &
    \dots &
    \dfrac{\partial \mu_1(0)}{\partial I_2(0)} &
    \dfrac{\partial \mu_1(0)}{\partial \gamma} &
    \dfrac{\partial \mu_1(0)}{\partial \theta} \\[10pt]
    \dfrac{\partial \mu_2(0)}{\partial S_1(0)} &
    \dots &
    \dfrac{\partial \mu_2(0)}{\partial I_2(0)} &
    \dfrac{\partial \mu_2(0)}{\partial \gamma} &
    \dfrac{\partial \mu_2(0)}{\partial \theta} \\[10pt]
    \dfrac{\partial \mu_1(1)}{\partial S_1(0)} &
    \dots &
    \dfrac{\partial \mu_1(1)}{\partial I_2(0)} &
    \dfrac{\partial \mu_1(1)}{\partial \gamma} &
    \dfrac{\partial \mu_1(1)}{\partial \theta} \\[10pt]
    \dfrac{\partial \mu_2(1)}{\partial S_1(0)} &
    \dots &
    \dfrac{\partial \mu_2(1)}{\partial I_2(0)} &
    \dfrac{\partial \mu_2(1)}{\partial \gamma} &
    \dfrac{\partial \mu_2(1)}{\partial \theta} \\[10pt]
    \vdots &  & \vdots & \vdots & \vdots
    \\[6pt]
    \dfrac{\partial \mu_2(T{-}1)}{\partial S_1(0)} &
    \dots &
    \dfrac{\partial \mu_2(T{-}1)}{\partial I_2(0)} &
    \dfrac{\partial \mu_2(T{-}1)}{\partial \gamma} &
    \dfrac{\partial \mu_2(T{-}1)}{\partial \theta}
  \end{pmatrix} \\
  &\text{ and } \\
  %% G_{(i,t),k} &= \partial \mu_i(t)/\partial \phi_k \\
  %
  %
  \mathbf{W} &= \text{diag} \left(\dots,\frac{\rho^2}{\sigma_i^2(t)} +                                                                      
  \frac{a_i^2(t)}{2\sigma_i^4(t)},\dots\right) \\
  %
  %
  %
  &= \text{diag}\left(\dots,
    \frac{\rho}{\mu_i(t) + \rho\mu_i^2(t)/k}
    + \frac{\bigl(1 + 2\rho\mu_i(t)/k\bigr)^2}
           {2\mu_i^2(t)\bigl(1 + \rho\mu_i(t)/k\bigr)^2},
           \dots\right).
  \end{split}
\end{align}
Note that $\mathbf{G}^T\mathbf{WG}$ is also the Fisher information
matrix for an ordinary least squares problem whose design matrix is
$\mathbf{W}^{1/2}\mathbf{G}$.

We denote the Fisher information for a single season $s$ as
$\mathbf{J}(s)$, so the corresponding Fisher information for $\theta$
is $1 / [\mathbf{J}^{-1}(s)]_{\theta\theta}$. When multiple seasons
are available, we assume independence across seasons, treating each
season's initial conditions as nuisance parameters. The Fisher
information for $\theta$ becomes:
\begin{equation}\label{eq:crlb_multi}
\text{I}(\hat{\theta}) = \sum_{s} \frac{1}{
  [\mathbf{J}^{-1}(s)]_{\theta\theta}},
\end{equation}
Seasons where the Fisher information computation fails (e.g., due to
near-zero incidence producing a singular Fisher information matrix)
contribute no information and are assigned
$[\mathbf{J}^{-1}(s)]_{\theta\theta} = \infty$.


%% Identifiability in systems-biology models is sometimes explored via
%% the profile likelihood~\citep{raue2009, shrestha2011} which is defined,
%% for a single season via
%% %
%% %
%% $$
%% \ell_{\text{PL}}(\theta) = \max_{S_1(0), I_1(0), S_2(0), I_2(0), \gamma}
%% \ell(S_1(0), I_1(0), S_2(0), I_2(0), \gamma, \theta).
%% $$
%% %
%% %
%% The profile likelihood is calculated numerically over a grid of
%% $\theta$ values, optimizing over all other parameters at each grid
%% point. Such a calculation would be computationally prohibitive at the
%% scale of our analysis. The CRLB is effectively a quadratic
%% approximation to the profile likelihood around the maximum-likelihood
%% estimate and is considerably cheaper to compute. Identifiability for
%% ODE-based epidemic models has also been studied via structural
%% identifiability~\citep{tuncer2018, chowell2023structural} and via
%% identifiable parameter combinations~\citep{eisenberg2014}; we use the
%% CRLB because it bounds the \emph{variance} of an estimator under a
%% hypothesized data design, which is the question we want to answer.

\subsection{Hypothesis testing for the SIR model}\label{subsec:hyp}
For each state pair and each of the two nested comparisons of
Table~\ref{tab:cases} we test $H_0: \theta = 0$ against $H_1: \theta >
0$. Under standard regularity conditions the maximum-likelihood
estimator $\hat\theta$ is asymptotically $\mathcal{N}(\theta,
I^{-1}(\theta))$.  We reject $H_0$ at level $\alpha = 0.05$ whenever
the one-sided $95\%$ Wald confidence interval $[\hat\theta -
  1.645 \cdot I^{-1/2}(\theta), \infty)$ excludes zero, i.e.
\begin{equation}\label{eq:wald}
  \hat\theta > z_{0.05}\cdot  I^{-1/2}(\theta).
\end{equation}
A likelihood ratio test using $\Lambda = 2\bigl[\ell(H_1) -
  \ell(H_0)\bigr]$ gives an asymptotically equivalent test.


\subsection{Simulated Analysis}\label{subsec:sim_crlb}
We simulate epidemics across a grid of coupling values $\theta \in
[10^{-4}, 10^{-1}]$ and seasonal forcing amplitudes $\varepsilon \in
[0, 1)$, comparing in-phase ($\varphi_1 = \varphi_2$) and
  reverse-phase ($\varphi_1 - \varphi_2 = \pi$) epidemics. At each
  grid point, the initial conditions $(S_i(0)/N_i, I_i(0)/N_i)$ for
  each simulated season and region are drawn independently from the
  empirical distribution of per-season, per-state initial conditions
  fitted to the real data (see SI~Fig.~\ref{fig:si_ics}). We repeat
  each configuration 500 times with fresh IC draws. We then compute
  optimal statistical power as described in Section
  \ref{subsec:hyp}. Note that $\theta$ is assumed known, so this
  analysis generates an optimistic estimate of statistical power. We
  report the mean power across the 500 IC realizations as a function
  of $(\theta, \varepsilon)$.

\subsection{Influenza Mortality Data}\label{subsec:real_data}
We analyze pneumonia and influenza mortality across 47 US states.  We
exclude Alaska and Hawaii since we do not have absolute humidity data
for them, and we exclude Arizona since it is a very arid state with
relatively inconsistent absolute humidity values (see Supplementary
Information). We integrate data from CDC FluView
(2010--2018)~\citep{cdcfluview2024} and HealthData.gov
(2017--2025)~\citep{healthdatagov2025}, excluding COVID pandemic years
2020--2022. Since pneumonia and influenza deaths include a year-round
non-influenza baseline, we estimate influenza burden as excess
mortality above a pre-season baseline, computed as the mean weekly
death count during July--October preceding each flu
season~\citep{serfling1963, viboud2006}. The infection fatality rate
$\rho \approx 5.36 \times 10^{-4}$ follows from the first row of Table
1 in~\citep{rolfes2018}. Seasonal phases $\varphi_i$ are estimated
from absolute humidity data~\citep{daon2025, shaman2009} by finding
the phase shift that maximizes the correlation between negative
absolute humidity values and a phase-shifted sine. We define each flu
season as $T = 20$ weekly observations starting
November~1$^{\text{st}}$ and fit a single coupling $\theta$ shared
across all seasons, with per-season initial conditions $(S_i(0),
I_i(0))$ inferred via maximum-likelihood. We aggregate the per-season
Cram\'er-Rao lower bounds via~\eqref{eq:crlb_multi} to obtain a bound
on the overall estimator variance for the coupling $\theta$.
 
\subsection{Maximum-Likelihood Estimation}\label{subsec:mle}
We maximize the log-likelihood~\eqref{eq:loglik} using
\texttt{SLSQP}~\citep{kraft1988} from
\texttt{NLopt}~\citep{NLopt}. Analytical gradients were calculated
via~\eqref{eq:grad}. Whenever an optimization resulted in
$\hat{\theta} \in [0.5,1]$, which is not physical, we simply folded
$\hat{\theta}$ via $\hat{\theta} \leftarrow 1-\hat{\theta}$. To handle
potential local minima, we use a multi-start strategy with 200 random
initial points and retain the solution with the highest
likelihood. Pilot experiments show that 20 initial points suffice, so
we take an order of magnitude more just to be safe.

\subsection{AI usage disclosure}
The author used Anthropic's Claude (Claude Code, Opus-4.7) to assist
with (i) drafting and refactoring analysis and simulation code (Python
and R), including the sensitivity-equation implementation and cluster
job scripts; (ii) drafting and copy-editing portions of the manuscript
and supplementary text; and (iii) formatting and figure layout in
\LaTeX. All model equations, derivations, statistical analyses,
figures, and interpretations were designed, verified, and approved by
the authors, who take full responsibility for the content of this
manuscript.

\section{Data and Code Availability}
Pneumonia and influenza mortality data are publicly available from CDC
FluView (2010--2018) and HealthData.gov (2017--2025); absolute
humidity data underlying phase estimation are
from~\citep{daon2025}. Code is available at the project's
\href{https://github.com/yairdaon/contacts}{repository}.

\section{Acknowledgments and Contributions}
The author thanks Michael Klots for early contributions to this
work. Computations were performed on the SLURM cluster at Bar-Ilan
University. Y.D. designed the research, performed the analysis, and
wrote the paper. This research was supported by The Israel Science Foundation
(grant No. 2332/26).


%% \section{Competing Interests}
%% The author declares no competing interest.

\bibliographystyle{plainnat}
\bibliography{~/lib.bib}

\end{document}

%% Lack of practical identifiability has previously been shown to hamper
%% reliable predictions in COVID epidemic
%% models~\citep{gallo2022}. 


%% \significancestatement{US public-health agencies forecast seasonal
%%   influenza using models that encode connections between states. We
%%   show that when states share a common seasonal driver --- as US
%%   influenza does --- the strength of inter-state connections is
%%   redundant once each state's epidemic starting conditions are
%%   accounted for. Two complementary analyses, one mechanistic and
%%   one model-free, agree that connections between states add no
%%   measurable skill to forecasts of US seasonal influenza. The same
%%   redundancy is likely in any system of coupled oscillators driven
%%   by a common force, from neural circuits to ecological
%%   metapopulations and climate networks. We conclude that
%%   inter-state coupling can be omitted from US influenza
%%   preparedness models without measurable loss of skill.}
